\documentclass[conference]{IEEEtran}
\IEEEoverridecommandlockouts
\usepackage{cite}
\usepackage{amsmath,amssymb,amsfonts}
\usepackage{textcomp}
\usepackage{booktabs}
\usepackage{url}
\usepackage{xcolor}

\hyphenation{op-tical net-works semi-conduc-tor}

\begin{document}

\title{Orbit-Transfer: Multi-Relay Bandwidth Aggregation\\ with Rateless Fountain Codes}

\author{\IEEEauthorblockN{Dylan Ondo}
\IEEEauthorblockA{\textit{Independent Researcher} \\
Libreville, Gabon \\
Email: danielebang59@gmail.com}}

\maketitle

\begin{abstract}
This paper presents Orbit-Transfer, a hybrid peer-to-peer and edge-relay file
transfer system that aggregates the capacity of multiple constrained network
paths using rateless fountain codes. A payload is split into $K$ source
symbols, augmented with $S$ LDPC precode checks, and the sender emits an
unbounded stream of encoded symbols; the receiver reconstructs the payload
from any $K+\varepsilon$ distinct symbols without retransmission requests.
Because fountain symbols are interchangeable and order-independent, the
sender can round-robin symbols across $N$ relays---and an optional direct
peer-to-peer channel---so the aggregate throughput is the sum of the live
paths. We evaluate the design on a multi-relay testbed in which each relay is
rate-limited by a token-bucket uplink shaper. Over five runs, two relays
deliver $2.2\times$ and three relays $3.2\times$ the single-relay throughput,
with negligible decoding overhead in the lossless systematic phase. We also
report a practical engineering finding: naive CPU-spinning rate limiters
collapse under concurrency, and we show how a high-resolution timer combined
with a self-correcting token bucket keeps independent relays paced uniformly.
\end{abstract}

\begin{IEEEkeywords}
erasure codes, fountain codes, LT codes, Raptor codes, bandwidth aggregation,
multipath transport, peer-to-peer, rate limiting.
\end{IEEEkeywords}

\section{Introduction}
\IEEEPARstart{L}{arge} transfers over constrained networks face two classic
obstacles: the bottleneck link between the endpoints and the retransmission
protocols that multiply latency and jitter. When a single path is
rate-limited---a residential uplink, a cellular quota, or a shared relay---the
obvious remedy is to use several of them at once. Multipath transport
(\texttt{MPTCP}~\cite{rfc6824}, SCTP) and resource pooling~\cite{wischik2008survey}
make this possible at the transport layer, but they require path coordination,
congestion coupling, and end-to-end transport state. In contrast, fountain
codes~\cite{byers1998digital,luby2002lt,mackay2005fountain} make aggregation
nearly trivial: symbols are interchangeable, so a sender may simply round-robin
them over every available path and let the receiver assemble the payload from
whatever arrives.

Orbit-Transfer exploits exactly this property. A sender splits a file into $K$
source blocks, applies a Raptor-style LDPC precode~\cite{shokrollahi2006raptor}
to obtain $L=K+S$ precoded inputs, and emits an \emph{unbounded} stream of
distinct symbols. The receiver needs only $K+\varepsilon$ \emph{any} symbols to
rebuild the file; it never asks for a missing block. The sender keeps emitting
until the receiver signals \textsc{Ready}, so any loss rate, reordering, or
path failure is absorbed automatically.

The practical consequence is a very simple multi-path scheduler: distribute
symbols round-robin across $N$ edge relays (and optionally a direct peer-to-peer
TCP channel), disable paths that fail, and rely on the rateless property to
finish. Because the relays do not need to coordinate, bandwidth aggregates
linearly: $N$ relays of capacity $C$ deliver up to $N\cdot C$, as long as each
relay is paced independently and accurately.

This paper makes the following contributions:
\begin{itemize}
\item A complete system description of Orbit-Transfer: a compact binary
protocol, a room-based edge relay with pre-receiver buffering and backpressure,
a multi-path scheduler, and per-symbol AEAD encryption (\S\ref{sec:arch}).
\item The fountain codec design---systematic phase, Robust Soliton
distribution, LDPC precoding, and peeling decoding---and its parameters
(\S\ref{sec:codec}).
\item An empirical aggregation study: two relays deliver $2.2\times$ and three
relays $3.2\times$ a single rate-limited relay, and the codec finishes with
essentially zero overhead in the systematic phase (\S\ref{sec:eval}).
\item A practical rate-limiting fix: CPU-spinning token buckets do not scale
under concurrency; replacing them with a high-resolution timer and a
self-correcting credit model keeps independent relays perfectly uniform
(\S\ref{sec:throttle}).
\end{itemize}

\section{System Architecture}
\label{sec:arch}

\begin{figure}[!t]
\centering
\scriptsize
\begin{minipage}{0.95\columnwidth}
\begin{verbatim}
Sender --s--> Relay 1 --s--> Receiver
  |              ...               ^
  +--s--> Relay N --s--------------+
  +--s--> direct (optional)--------+
  <---------- READY ---------------+
\end{verbatim}
\end{minipage}
\caption{System model. The sender distributes symbols over $N$ relays and an
optional direct peer-to-peer channel; the receiver merges all feeds and signals
\textsc{Ready} once decoded.}
\label{fig:system}
\end{figure}

The system (Fig.~\ref{fig:system}) comprises three roles: a \emph{sender}, one
or more \emph{relays}, and a \emph{receiver}. The sender connects to every
relay over WebSocket and, when a direct path is advertised, to the receiver
over raw TCP. Both sides share a 64-bit session identifier.

\subsection{Protocol}
Binary frames are carried over WebSocket (relay links) or raw TCP (direct
P2P), with the layout
\texttt{[type: u8][len: u32 LE][payload]}. Table~\ref{tab:protocol} lists the
message types. The \textsc{Meta} frame carries the filename, payload size,
symbol size, $K$, and an XXH3 checksum; \textsc{Symbol} carries an encoding
symbol; \textsc{Ready} terminates the transfer; \textsc{Direct} advertises the
receiver's peer-to-peer listener address; \textsc{Done} tears a session down.

\begin{table}[!t]
\centering
\caption{Protocol messages.}
\label{tab:protocol}
\begin{tabular}{@{}lp{0.62\columnwidth}@{}}
\toprule
Message & Payload \\
\midrule
\textsc{Hello} & session id (u64), role (sender/receiver) \\
\textsc{Meta} & session id, filename, size, symbol size, $K$, checksum \\
\textsc{Symbol} & session id, esi (u32), symbol data \\
\textsc{Ready} & session id \\
\textsc{Direct} & session id, advertised address \\
\textsc{Done} & session id \\
\bottomrule
\end{tabular}
\end{table}

\subsection{Relay}
Each relay maintains a map of \emph{rooms} keyed by session id, with one
sender link and one receiver link per room. Because the receiver may join
later, the relay buffers symbols in a bounded queue ($\le 256$\,MiB) and flushes
them in order on registration. Routing is entirely by session id; the relay is
stateless with respect to the payload and never decodes it. Per-connection
outgoing channels provide backpressure that propagates to TCP flow control, so
a slow receiver or a rate-limited link never causes silent drops---symbols
queue and the sender is eventually throttled rather than losing data.

The relay also exposes a per-connection egress rate limiter
(\S\ref{sec:throttle}) used in our experiments to emulate a constrained uplink.

\subsection{Sender and receiver sessions}
The sender session attaches to all relays, opens per-path bounded queues
(1024 symbols) with one writer task each, emits \textsc{Meta}, and then streams
symbols. Each writer task drains its queue independently at the path's own
rate, so the aggregate is the sum of live paths; a failed send marks the path
dead and its load is redistributed automatically.

The receiver session spawns one reader task per relay that re-emits frames on
a single channel, so a \texttt{select}-based decode loop merges all paths (and
the optional direct channel) into one ordered feed. The direct path runs over
raw TCP by default and over \textbf{QUIC} when the receiver advertises a
\texttt{quic://} address (the sender picks the transport from the prefix);
both carry the same binary framing. After decoding, it sends \textsc{Ready} on
every link and then \emph{drains} each socket to EOF before closing---closing
with unread data would emit a TCP reset and purge the in-flight
\textsc{Ready}.

\section{Rateless Fountain Codec}
\label{sec:codec}

\subsection{Encoding}
The payload is split into $K$ source symbols of equal size. A Raptor-style
LDPC precode~\cite{shokrollahi2006raptor,rfc5053} computes
\begin{equation}
S = \max\!\bigl(2,\ \lfloor K/100\rfloor + 10\bigr)
\end{equation}
check symbols; check $i$ is the XOR of the source symbols whose index is
congruent to $i$ modulo $S$ together with sources $i\bmod K$ and $(i{+}1)\bmod K$.
A RaptorQ-style HDPC~\cite{rfc6330} layer then appends $H$ high-density checks,
where $H$ is the smallest integer such that
\begin{equation}
\binom{H}{\lfloor H/2\rfloor} \ge K + S,
\end{equation}
and check $i$ XORs a dense pseudo-random subset (about half) of the $K+S$
precoded inputs, seeded deterministically from the check index. The resulting
$L=K+S+H$ precoded symbols are the inputs to an LT encoder~\cite{luby2002lt}
that draws its degree from the Robust Soliton Distribution with
$c=0.1$, $\delta=0.05$, and $R=c\sqrt{K}\ln(K/\delta)$. The HDPC checks act as
the decoding graph's safety net: at small $K$ they let reconstruction finish
near $K+\varepsilon$ instead of stalling on the sparser LDPC graph.

The code is \emph{systematic}: for esi $<K$ the encoder emits the raw source
symbol directly, so in the lossless case the receiver reconstructs the payload
from the $K$ systematic symbols with \emph{zero} overhead. For esi $\ge K$ it
emits random combinations. Symbol generation is deterministic and
reproducible: the neighbor selection is seeded from the esi with a
ChaCha8 PRNG, so any party can recompute the encoding graph from an esi
without a shared dictionary.

\subsection{Decoding}
The decoder builds the decoding graph over the $L$ precoded inputs. It seeds
the peeling process with the $S$ LDPC and $H$ HDPC check equations and iterates
degree-one equations: resolving a symbol XORs it out of every equation that
mentions it, and a degree-one equation resolves immediately. The dense HDPC
equations resolve last, guaranteeing the graph completes. Reconstruction
completes once all $K+S+H$ inputs are resolved; the payload is then verified
against the XXH3 checksum carried in \textsc{Meta}. Duplicate and out-of-order
symbols are harmless (they are de-duplicated by esi), which is the property
that makes multi-path aggregation trivial. The core XOR kernel is
SIMD-vectorizable, and the encode rate exceeds $3$\,Gb/s on commodity hardware
with a 64\,MiB payload.

\section{Multi-Path Aggregation}
\label{sec:agg}

Aggregation in Orbit-Transfer reduces to a scheduling decision, because
fountain symbols are interchangeable and order-independent. The sender uses a
deterministic scheduler~\cite{maymounkov2002online} that round-robins the relay
load across the $N$ relays while keeping a relay/direct ratio. Two properties
matter in practice:
\begin{itemize}
\item \textbf{Dead-path fallback.} A relay that dies mid-transfer is disabled;
its queued symbols are lost, but the receiver never notices---it needs any
$K+\varepsilon$ symbols, not a specific subset. The scheduler simply moves on.
\item \textbf{Drain-based shutdown.} After \textsc{Ready}, the receiver drains
every relay link to EOF so relays flush their queued symbols cleanly rather
than hitting a TCP reset.
\end{itemize}

Because each relay is paced independently, the expected aggregate is $N\cdot C$.
The experimental section (\S\ref{sec:eval}) confirms this holds to within the
variance of the per-relay limiters.

\section{Relay Rate Limiting}
\label{sec:throttle}

Our experiments emulate constrained uplinks with a per-connection token-bucket
shaper at the relay: each symbol costs its byte count, and the bucket refills
at the configured rate. Two implementation details turned out to be critical.

\subsection{Spin-yield pacing does not scale}
The first implementation paced a symbol by spinning on
\texttt{Instant::now()} while yielding to the scheduler:
\begin{verbatim}
while Instant::now() < deadline {
    yield_now().await;
}
\end{verbatim}
Alone, one relay paced at $\approx 343$ symbols/s instead of the nominal
$500$/s. Under three concurrent relays the per-relay rates collapsed to
$361/316/350$ symbols/s ($1028$/s total instead of $1500$/s): each spinning task
contends for CPU, and the scheduling overhead grows with the number of
concurrent spinners. The limiter both undershoots and loses independence, which
directly undermines the aggregation claim.

\subsection{Parked sleeps with a self-correcting bucket}
We replaced the spin with a real \texttt{sleep} and raised the Windows timer
resolution to 1\,ms via \texttt{timeBeginPeriod(1)}. Sleeping parks the task,
so concurrent relays no longer contend for CPU. Yet the \texttt{tokio} timer
still sleeps $\approx 3.7$\,ms for a requested $2$\,ms, which alone would
undershoot the nominal rate by half.

The fix is a token bucket whose accounting is \emph{self-correcting}: the
credit window is closed at the \emph{start} of every call and \emph{never}
after the sleep, so whatever time the sleep actually took is credited back to
the next symbol. The long-term rate therefore converges to the configured
value regardless of the timer's granularity. A burst cap ($\approx 10$\,ms of
credit) prevents a long idle from unlocking an unbounded run. The result is a
uniform $613/613/614$ symbols/s across three concurrent relays---a perfect
$3\times$ scale-up with no contention.

\section{Experimental Evaluation}
\label{sec:eval}

\subsection{Setup}
All experiments run on a single Windows 11 machine over loopback, in release
mode (debug builds slow the encoder by an order of magnitude). Each relay
rate-limits its egress to a nominal $2$\,MiB/s. The payload is 16\,MiB with
4096-byte symbols ($K=4096$). We report five runs of the aggregation bench;
the table gives the mean and range. Aggregation factors are computed per-run
as the ratio of elapsed times (payload size is identical across cases).

\subsection{Codec overhead}
In the lossless case the systematic phase alone reconstructs the payload:
the decoder finishes at exactly $K$ symbols (overhead $1.000$). Under simulated
$5\%$ loss the overhead is $\approx 1.1\times L$, well inside the Raptor
guarantee. Encode throughput exceeds $3$\,Gb/s for a 64\,MiB payload; decode of
the same payload completes in $\approx 0.2$\,s.

\subsection{Aggregation}
\begin{table}[!t]
\centering
\caption{Throughput vs.\ number of rate-limited relays (16\,MiB payload,
4096-B symbols, mean of 5 runs).}
\label{tab:bench}
\begin{tabular}{@{}lcccc@{}}
\toprule
Case & MiB/s & Range & Symbols & Factor \\
\midrule
1 relay & 2.29 & 2.19--2.34 & $\approx$21.6k & 1.00 \\
2 relays & 5.01 & 4.31--5.39 & $\approx$34.5k & 2.19 \\
3 relays & 7.24 & 6.66--7.78 & $\approx$34.6k & 3.16 \\
3 relays + direct & 5.76 & 5.40--6.30 & $\approx$43.9k & 2.51 \\
\bottomrule
\end{tabular}
\end{table}

Table~\ref{tab:bench} reports the results. Two relays deliver $2.19\times$
(range $1.84$--$2.47\times$) and three relays $3.16\times$ (range
$2.84$--$3.70\times$) the single-relay baseline, matching the expected $N\times$
scaling to within the per-relay limiter variance. The single-relay case is
remarkably stable ($2.19$--$2.34$\,MiB/s across runs).

The ``3 relays + direct'' case adds the receiver's peer-to-peer listener; the
direct channel carries only the redundant tail (esi $\ge K$), because the
receiver completes at $\approx K$ symbols and the throttled relays already
saturate the decode feed. Consequently direct P2P adds little on this loopback
testbed. It is intended for the opposite scenario---when relays are the
unavailable or the bottleneck path, direct P2P is the faster lane with
automatic fallback---rather than as a multiplier on top of saturated relays.

\subsection{Rate-limiter independence}
Table~\ref{tab:throttle} isolates the limiter by pacing three relays in
parallel with no sessions or decoding. The per-relay rates are within $0.4\%$
of each other and sum to $3\times$ the single-relay rate, confirming that the
aggregation is limited only by the per-path pipes and not by any shared CPU
resource.

\begin{table}[!t]
\centering
\caption{Per-relay egress with three concurrently throttled relays.}
\label{tab:throttle}
\begin{tabular}{@{}lcc@{}}
\toprule
Relays & Per-relay rate & Total \\
\midrule
1 & 624 sym/s & 624/s \\
3 & 613, 613, 614 sym/s & 1841/s \\
\bottomrule
\end{tabular}
\end{table}

\section{Related Work}
Fountain codes originate with the digital-fountain framework of
Byers \emph{et al.}~\cite{byers1998digital} and the LT construction of
Luby~\cite{luby2002lt}; MacKay gives an accessible survey~\cite{mackay2005fountain}.
Raptor~\cite{shokrollahi2006raptor,rfc5053} and RaptorQ~\cite{rfc6330} codes add
a fixed precode to reduce decoding overhead to a few percent and are used in
3GPP broadcast and reliable delivery standards. Our codec follows the Raptor
recipe (LT over an LDPC precode, systematic phase) at smaller scale and with a
deterministic, reproducible neighbor seed.

Multipath transport has been standardized as
\texttt{MPTCP}~\cite{rfc6824} and studied through the resource-pooling lens of
Wischik \emph{et al.}~\cite{wischik2008survey}. In contrast to transport-layer
coupling, our scheme needs no congestion coordination between paths: the
rateless code absorbs the path-to-path rate differences, and the only shared
state is the session id.

Fountain codes have long been proposed for peer-to-peer bulk
delivery~\cite{maymounkov2002online}; Orbit-Transfer makes the multi-path
extension explicit and quantifies it on rate-limited relays, which to our
knowledge is not reported in the fountain literature.

\section{Conclusion and Future Work}
Orbit-Transfer shows that rateless fountain codes turn multi-path aggregation
into a scheduling problem with an essentially linear payoff: $N$ rate-limited
relays deliver $\approx N\times$ the throughput of one, because symbols are
interchangeable and each path is paced independently. The engineering
takeaways are (i) the systematic Raptor-style precode removes overhead entirely
in the lossless case, and (ii) rate limiters must be parked-sleep based and
self-correcting to keep independent paths uniform under concurrency.

Ongoing and future work includes extending the direct path beyond TCP and QUIC
(already implemented: the receiver may advertise a \texttt{quic://} listener and
the sender picks QUIC automatically, giving 0-RTT setup and connection
migration), and pushing the RaptorQ-style precode toward even lower overhead at
very small $K$. The relay is already transport-agnostic, and the protocol is
compact enough to carry over any ordered or unordered channel.

In addition to the native multi-path stack, a WebRTC bridge is shipped with the
WebAssembly package: browser-to-browser transfers run over an SCTP data channel
with the same fountain codec compiled to WASM. Signaling is automated via a
tiny in-memory room on the static server: the sender drags a file and gets a
shareable \texttt{?room=<id>} link, the receiver opens it, and the SDP
offer/answer is exchanged over HTTP. A practical lesson emerged from browser
testing: Chrome masks host ICE candidates behind mDNS names
(\texttt{*.local}) that some network interfaces cannot resolve, leaving the
ICE agent with no candidate pair (stuck in the ``new'' state); since candidate
ports are not masked, adding an explicit loopback candidate pointing at the
peer's port restores connectivity between tabs on the same host. The bridge is
validated end-to-end in the browser and through mocked-channel integration
tests (including the automated share-link flow).

\begin{thebibliography}{1}
\bibitem{byers1998digital}
J.~W. Byers, M.~Luby, M.~Mitzenmacher, and A.~Rege, ``A digital fountain approach
to reliable distribution of bulk data,'' in \emph{Proc. ACM SIGCOMM}, 1998,
pp.~56--67.
\bibitem{luby2002lt}
M.~Luby, ``LT codes,'' in \emph{Proc. 43rd IEEE Symp. Foundations of Computer
Science (FOCS)}, 2002, pp.~271--280.
\bibitem{mackay2005fountain}
D.~J.~C. MacKay, ``Fountain codes,'' \emph{IEE Proc. Communications},
vol.~152, no.~6, pp.~1062--1068, 2005.
\bibitem{shokrollahi2006raptor}
A.~Shokrollahi, ``Raptor codes,'' \emph{IEEE Trans. Information Theory},
vol.~52, no.~6, pp.~2551--2567, 2006.
\bibitem{rfc5053}
M.~Luby, A.~Shokrollahi, M.~Watson, and T.~Stockhammer, ``Raptor forward error
correction scheme for object delivery,'' IETF RFC 5053, 2007.
\bibitem{rfc6330}
M.~Luby, A.~Shokrollahi, M.~Watson, and T.~Stockhammer, ``RaptorQ forward error
correction scheme for object delivery,'' IETF RFC 6330, 2011.
\bibitem{rfc6824}
A.~Ford, C.~Raiciu, M.~Handley, and O.~Bonaventure, ``TCP extensions for
multipath operation with multiple addresses,'' IETF RFC 6824, 2013.
\bibitem{wischik2008survey}
D.~Wischik, M.~Handley, and M.~Bagnulo Braun, ``The resource pooling
principle,'' \emph{ACM SIGCOMM Comput. Commun. Rev.}, vol.~38, no.~5,
pp.~47--52, 2008.
\bibitem{maymounkov2002online}
P.~Maymounkov and D.~Mazi{\`e}res, ``Rateless codes and big downloads,'' in
\emph{Proc. 2nd Int. Workshop on Peer-to-Peer Systems (IPTPS)}, 2003,
pp.~247--255.
\bibitem{webrtc}
A.~Bergkvist \emph{et al.}, ``WebRTC 1.0: real-time communication between
browsers,'' W3C Recommendation, 2021.
\end{thebibliography}

\end{document}